Nuclear medicine exploits the physiological accumulation of distinct radiopharmaceuticals to functionally evaluate specific organs by utilizing internally administered unsealed sources of radiation. In the context of thyroid evaluation, the primary mechanism of interest is the biological trapping and subsequent organification of iodine by thyroid follicular cells. Iodine is a fundamental building block for the synthesis of thyroid hormones (triiodothyronine [$\text{T}_3$] and thyroxine [$\text{T}_4$]). By administering known activities of radioactive iodine isotopes—typically $^{123}\text{I}$ or $^{131}\text{I}$ (administered activity ranging from 10 $\mu$Ci to 30 $\mu$Ci)—the fractional uptake of the tracer by the thyroid gland can be quantitatively analyzed over specific temporal intervals (e.g., 2, 6, and 24 hours post-administration). This fraction, commonly referred to as the Radioactive Iodine Uptake (RAIU), provides profound quantitative insight into hyperthyroidism, hypothyroidism, and the pathophysiological status of nodular anomalies \cite{cherry2012physics, captus4000e}.

\subsection{Physical Properties of Radionuclides Used}
Two principle isotopes of iodine are routinely utilized for uptake measurements:
\begin{itemize}
    \item \textbf{$^{131}\text{I}$:} A beta-minus ($\beta^-$) and gamma ($\gamma$) emitter with a physical half-life ($T_{1/2}$) of 8.02 days. The principal gamma emission utilized for counting and imaging is 364 keV (81\% abundance). Because of its significant particulate emission, its primary role in contemporary nuclear medicine is therapeutic (e.g., ablation of hyperthyroid tissue or thyroid carcinoma post-thyroidectomy), though it is historically and fundamentally utilized in uptake studies.
    \item \textbf{$^{123}\text{I}$:} An electron capture (EC) decaying nuclide with a physical half-life of 13.2 hours. It decays to $^{123}\text{Te}$, emitting a primary gamma ray of 159 keV (83\% abundance). Its absence of particulate radiation translates to a remarkably lower patient radiation dosimetry compared to $^{131}\text{I}$, making it the preferred tracer for purely diagnostic imaging and uptake measurements.
\end{itemize}

\subsubsection{Cesium-137 ($^{137}$Cs)}
$^{137}$Cs is a standard calibration source extensively used in nuclear instrumentation. It undergoes $\beta^-$ decay to metastable $^{137m}$Ba, which subsequently emits a gamma photon of 662 keV.
Key characteristics include:
\begin{itemize}
    \item Half-life: approximately 30.1 years
    \item Main gamma energy: 662 keV
\end{itemize}
The monoenergetic gamma emission and long half-life make $^{137}$Cs highly suitable for energy calibration, resolution measurement, and constancy checks of scintillation detectors.

\subsubsection{Europium-152 ($^{152}$Eu)}
$^{152}$Eu is a multi-energy gamma emitter widely used for linearity and energy calibration of detection systems. It decays via $\beta^-$ and electron capture processes with a half-life of approximately 13.5 years.
It emits multiple gamma rays across a broad energy spectrum, including:
\begin{itemize}
    \item 121 keV, 244 keV, 344 keV, 778 keV, 964 keV, 1408 keV (among others)
\end{itemize}
The presence of multiple well-defined gamma peaks makes $^{152}$Eu particularly useful for assessing the linearity of the multichannel analyzer and verifying the accuracy of energy calibration over a wide range.

\subsection{Role of Radionuclides in Quality Assurance}
Different radionuclides serve distinct purposes in the quality assurance program of a thyroid uptake system:
\begin{itemize}
    \item $^{137}$Cs is primarily used for energy calibration, resolution measurement, and constancy testing due to its stable and well-defined gamma emission.
    \item $^{152}$Eu is used for linearity testing because of its multiple gamma energies spanning a wide range.
    \item $^{125}$I is useful in low-energy calibration and well detector performance evaluation.
    \item $^{131}$I is used in actual thyroid uptake measurements due to its physiological uptake by the thyroid gland.
\end{itemize}
The combined use of these radionuclides ensures accurate calibration, reliable detector performance, and precise quantification of thyroid uptake measurements.

\begin{figure}[H]
    \centering
    \includegraphics[width=0.6\textwidth]{"/Users/souvik/Desktop/Lab-Reports/4th Sem/Thyroid uptake/Figures/Theory/Sources.png"}
    \caption{Sources used for calibration and quality assurance.}
    \label{fig:sources_qa}
\end{figure}

\subsection{Thyroid Probe}
The thyroid probe is a specialized radiation counter predominantly employed to measure the fractional uptake of radioiodine, such as $^{131}\text{I}$ or $^{123}\text{I}$, within the thyroid gland following the oral administration of a radioactive capsule (e.g., $^{131}\text{I}$-NaI or $^{123}\text{I}$-NaI) \cite{saha_nuclear_medicine}. 

Fundamentally, the operational architecture of a thyroid probe largely mirrors that of a well counter. It incorporates a Thallium-activated Sodium Iodide, NaI(Tl), scintillation detector, typically possessing dimensions of $5\text{ cm}$ in diameter and $5\text{ cm}$ in thickness. This detector is optically coupled to a Photomultiplier (PM) tube. The associated electronic processing chain comprises a pre-amplifier, an amplifier, a Pulse Height Analyzer (PHA), and a scaler to register the quantitative counts \cite{saha_nuclear_medicine}.

One of the critical distinguishing features of the thyroid probe, which differentiates it from a standard well counter, is the necessity of a collimator. The collimator restricts the field of view specifically to the localized region of the thyroid gland. It structurally consists of a cylindrical lead barrel, approximately $20$ to $25\text{ cm}$ in length, which envelops both the NaI(Tl) crystal and the PM tube. This robust shielding significantly mitigates the background activity arising from $\gamma$-radiations originating from physiological areas outside the thyroid gland from reaching the detector \cite{saha_nuclear_medicine}.

The geometric detection efficiency of a thyroid probe operates inversely with the square of the distance separating the face of the detector and the patient's thyroid. During preparation, the probe is initially calibrated for photon energies utilizing a known standard, typically the $662\text{ keV}$ $\gamma$-ray peak of a $^{137}\text{Cs}$ source. Following baseline calibration, the system's discriminator settings (the PHA window) are specifically tuned to the characteristic energy of the administered isotope, such as the $364\text{ keV}$ $\gamma$-ray for $^{131}\text{I}$. 

While unscattered primary photons are targeted, photons that undergo Compton scattering within the thyroid tissue itself may still reach and interact within the detector because they originate within the unshielded field of view and are not stopped by the collimator barrel. To prevent these lower-energy scattered photons from falsely elevating the count rate, they are systematically excluded from the total measured counts by selecting highly appropriate lower and upper discriminator settings on the PHA carefully bounding the $364\text{ keV}$ $\gamma$-ray photopeak of $^{131}\text{I}$ \cite{saha_nuclear_medicine}. Additionally, the inherent attenuation of outgoing $\gamma$-photons by the overlying thyroid tissues themselves inherently reduces the overall detection efficiency of the internal activity.

\begin{figure}
    \centering
    \includegraphics[width=0.8\textwidth]{/Users/souvik/Desktop/Lab-Reports/4th Sem/Thyroid uptake/Figures/Theory/TyroidProbe.png}
    \caption{Schematic representation of a thyroid probe system, illustrating the NaI(Tl) scintillation crystal, photomultiplier tube, lead collimator, and associated electronic processing units. The collimator restricts the field of view to the thyroid region, while the scintillation crystal converts incident gamma photons into visible light, which is then amplified by the PMT for quantitative analysis.}
    \label{fig:thyroid_probe}
\end{figure}
\subsection{Apparatus and System Architecture}
The precise derivation of fractional radioiodine uptake requires highly specialized radiometric instrumentation. A typical Thyroid Uptake Measurement Unit (such as the CAPTUS 4000e unit) encompasses a meticulously engineered combination of scintillation detection, complex electronics, geometric shielding, and digital processing units. The overall architectural pipeline maps incident radiation into quantitative count rates within specific energy windows.

\begin{figure}[H]
    \centering
    \includegraphics[width=0.6\textwidth]{"/Users/souvik/Desktop/Lab-Reports/4th Sem/Thyroid uptake/Figures/Theory/CAPTUS 4000e Unit.png"}
    \caption{The CAPTUS 4000e Thyroid Uptake System.}
    \label{fig:captus4000e}
\end{figure}

The detailed technical specifications of the CAPTUS 4000e system are provided in Table \ref{tab:captus_specs}.

\begin{table}[H]
    \centering
    \renewcommand{\arraystretch}{1.3}
    \caption{Technical Specifications of the CAPTUS 4000e Unit}
    \label{tab:captus_specs}
    \begin{tabular}{|l|l|l|l|}
        \hline
        \multicolumn{3}{|l|}{\textbf{Specification}} & \textbf{Detailed} \\ \hline
        \multicolumn{3}{|l|}{Make} & Mirion - Capintec \\ \hline
        \multicolumn{3}{|l|}{Model} & CAPTUS 4000e \\ \hline
        Detector & Probe & Detector Material & NaI(Tl) \\ \cline{3-4}
         &  & Dimension & 2"X2" \\ \cline{3-4}
         &  & Resolution & $\leq$9.5\% or better for Cs137 \\ \cline{2-4}
         & Well & Detector Material & NaI(Tl) well-type crystal \\ \cline{3-4}
         &  & Dimension & 2"X2" \\ \cline{3-4}
         &  & WELL Diameter & 1.6CM \\ \cline{3-4}
         &  & WELL Depth & 3.8CM \\ \cline{3-4}
         &  & Resolution & $\leq$10.0\% or better for Cs137 \\ \cline{3-4}
         &  & Shielding Material & Lead lined with brass to minimise backscattering \\ \hline
        MCA & \multicolumn{2}{l|}{Type} & CAP MCA \\ \cline{2-4}
         & \multicolumn{2}{l|}{Channel} & 1024 \\ \cline{2-4}
         & \multicolumn{2}{l|}{Maximum Count Rate} & 200000Cps \\ \cline{2-4}
         & \multicolumn{2}{l|}{ROI's Selection} & Automatic and Manual \\ \hline
        Operational & \multicolumn{2}{l|}{Temperature Range} & 10$^\circ$C to 30$^\circ$C \\ \cline{2-4}
        Environment & \multicolumn{2}{l|}{Pressure Range} & 91 to 105 KiloPascal \\ \cline{2-4}
        Requirement & \multicolumn{2}{l|}{Humidity Range} & Maximum relative humidity $<$ 95\% \\ \hline
    \end{tabular}
\end{table}

\subsubsection{Thallium-Activated Sodium Iodide Scintillator [NaI(Tl)]}
At the core of the measurement probe is the inorganic scintillation crystal, commonly a $2^{\prime\prime} \times 2^{\prime\prime}$ right cylindrical Thallium-activated Sodium Iodide [NaI(Tl)] detector. The intrinsic physical property of a scintillator is its capacity to absorb high-energy incident ionizing radiation (e.g., 364 keV photons) and iteratively dissipate that energy via the emission of secondary photons within the visible electromagnetic spectrum. 
\par
Incident gamma photons enter the crystal lattice and undergo two primary interactions:
\begin{enumerate}
    \item \textbf{Photoelectric Effect:} Dominant for lower energy photons ($< 200$ keV, particularly pertinent for $^{123}\text{I}$'s 159 keV emission), where the photon completely transfers its kinetic energy to an inner-shell electron (usually the K-shell) of an Iodine atom ($Z=53$). The ejected photoelectron propagates a Coulombic cascade through the crystal.
    \item \textbf{Compton Scattering:} More probable for intermediate energies (such as the 364 keV of $^{131}\text{I}$ and 662 keV calibration photons of $^{137}\text{Cs}$). The incident photon scatters elastically against a loosely bound outer-shell electron, depositing only a fraction of its energy. The scattered photon may subsequently undergo photoelectric absorption or escape the crystal boundary, yielding a partial energy deposition event (the "Compton continuum") \cite{knoll2010radiation}.
\end{enumerate}
The kinetic energy carried by the resulting high-speed electrons creates thousands of electron-hole pairs across the valence and conduction bands of the NaI lattice. Pure NaI exhibits a wide bandgap and emits light inefficiently and outside the ideal visible range; thus, Thallium (Tl) impurities are introduced strictly as luminescence activation centers. Electrons migrate through the conduction band into the Thallium traps and de-excite, emitting luminescence photons predominantly in the blue region (peak wavelength $\approx 415$ nm). The total number of visible photons successfully generated is directly and strictly proportional to the energy deposited by the initial gammas, solidifying the crystal's role as a potent energy spectrometer \cite{cherry2012physics}.

\subsubsection{Photomultiplier Tube (PMT) Dynamics}
The faint, instantaneous bursts of 415 nm scintillation light are optically coupled into a robust Photomultiplier Tube (PMT). The face of the PMT consists of a thin bialkali photo-emissive layer (the photocathode). Upon absorbing the scintillation photons, electrons are liberated via the photoelectric effect (now generating visible-light photoelectrons). Because the photo-conversion yield is small, these weak electron clusters are subjected to extreme electro-multiplication through a sequential array of dynodes. 
\par
A high negative voltage bias (typically 800V to 1200V) strictly orchestrates an accelerating electric field longitudinally down the tube. Electrons aggressively strike the first dynode, executing secondary electron emission, thereby tripling or quadrupling the localized charge. This avalanche iterates successively across 10-12 dynodes forming an exponential magnitude multiplication (often exceeding $10^6$ gain). The robust current pulse safely collected at the final anode represents an analog voltage pulse whose amplitude unequivocally mirrors the absolute value of the primary gamma photon's transferred energy \cite{knoll2010radiation, cherry2012physics}.

\subsubsection{Flat-Field Collimation and Shielding Criteria}
In stark contrast to multi-hole collimators seen entirely in gamma cameras utilized for spatial mapping, a thyroid uptake probe universally employs a \textbf{Flat-Field Collimator}. Its solitary function is generating an effectively uniform geometric counting efficiency over a macroscopic region (corresponding anatomically to the anterior neck and the spatially distributed thyroid lobes). 
\par
The probe housing is encased in substantive dense Lead ($\text{Pb}$) shielding to aggressively suppress external background interference (e.g., cosmic gammas or extraneous scatter generated from the room). A crucial engineering aspect is that this lead shielding is internally \textit{lined with brass or copper}. When high-energy gammas undergo photoelectric absorption in the primary Lead shielding, 72 to 88 keV characteristic X-rays are ejected. Brass possesses a drastically lower atomic number ($Z$) and effectively absorbs these secondary Lead X-rays, preventing them from corrupting the spectral baseline of the system detector.

\subsubsection{Neck Phantoms and Measurement Geometry}
Quantitative standardized measurements inherently require accounting precisely for tissue attenuation and systemic geometric spread. Standard Lucite neck phantoms (e.g., Biodex Neck Phantoms) structurally synthesize the physical photon attenuation parameters strictly characteristic of human neck tissue structurally mirroring an embedded vocal/tracheal tract. Standard baseline counts are acquired placing a precise calibration source (or identically administered patient capsule formulation) at a universally standardized depth within the lucite phantom. 

\begin{figure}[H]
    \centering
    \includegraphics[width=0.6\textwidth]{"/Users/souvik/Desktop/Lab-Reports/4th Sem/Thyroid uptake/Figures/Theory/Neck phantom.png"}
    \caption{A typical Lucite neck phantom used for standardizing measurement geometry and tissue attenuation.}
    \label{fig:neck_phantom}
\end{figure}
\par
The distance from the anterior surface of the target source to the face of the crystal collimator is strictly fixed, usually at \textbf{25 cm (or 30 cm)} utilizing a fixed spatial pointer rod. The inverse-square law dictates that minor positional translations produce exponentially scaled variances in observed flux rate; thus, anchoring the physical geometry identically across standard-dose measurement and subsequent chronological in-vivo patient measurements is mathematically crucial for RAIU consistency \cite{captus4000e}.

\subsubsection{The Multi-Channel Analyzer (MCA)}
The output amplified signals from the PMT are processed through a pre-amplifier and shaping amplifier, ultimately feeding into a Multi-Channel Analyzer. The MCA functions essentially as a high-speed analog-to-digital converter (ADC) combined with continuous memory registers. An MCA typically divides real-time pulse-height amplitudes into $1024$ distinct voltage segments or "channels" (equivalent roughly to energies spanning 0 keV to 2000 keV). As pulses arrive consecutively, the MCA tallies their occurrences into their corresponding channel width. This constructs an absolute energy spectrum—allowing operators (and internal analytic software) to strictly impose a Region of Interest (ROI) window around precisely desired photopeaks (e.g., spanning exactly 320 to 400 keV to singularly capture $^{131}\text{I}$ emissions), functionally severing all background Compton debris out of calculation consideration.

\subsection{System Calibration and Continuous Quality Assurance Standards}
To maintain the requisite regulatory precision alongside high clinical diagnostic yield, systematic procedures act as rigorous QA checks \cite{captus4000e}. The probe performance inevitably drifts functionally due to minute fluctuations in PMT high-voltage bias, ambient room temperatures modifying crystal transparency, and gradual hygroscopic NaI(Tl) degradation.

\begin{itemize}
    \item \textbf{Auto-calibration \& Energy Resolution Validation:} Utilizing a robust mono-energetic standard source (specifically the 662 keV gamma from $^{137}\text{Cs}$ having a heavily reliable $T_{1/2} = 30.08$ years). The probe processing module continually iterates software gain parameters and signal-zero offsets until the distinct 662 keV emission falls identically on a specified operational channel number (e.g., channel 331 indicating a resolution mapping of roughly 2 keV/channel). Beyond centering, the MCA functionally extracts the Full Width at Half Maximum (FWHM) of that photopeak. Resolution percent is formulated mathematically as:
    \begin{equation}
        \% \text{ Resolution} = \left( \frac{\text{FWHM}}{E_\gamma} \right) \times 100
    \end{equation}
    Under rigorous factory calibration specifications, optimal fractional resolution for a planar $2^{\prime\prime}\times2^{\prime\prime}$ NaI(Tl) probe must rigidly trace under $\leq 9.5\%$. Significant spreading signifies structural degradation such as crystal yellowing, fracture, or hydration.

    \item \textbf{Linearity Tests:} Confirms that the linear proportionality inherently established by the calibration extends equally onto lower and higher energy channels across the 1024 spectrum. It uses multi-peak reference standard sources such as $^{152}\text{Eu}$ or sequentially interrogating $^{133}\text{Ba}$ / $^{57}\text{Co}$. Substantial deviation unmasks severe PMT non-linearity.

    \item \textbf{Constancy Stability Verification:} Daily qualitative verification to preemptively reject equipment drift. It algorithmically estimates the current decayed true activity of a registered standard control source (usually $^{137}\text{Cs}$) from its initially registered value matrix. Measuring the source establishes a recorded activity utilizing embedded efficiency tables. Deviation constraints dictate that variations inherently cannot drift beyond $\pm 10\%$ of standard bounds over multi-day operations.

    \item \textbf{Poisson Statistical Congruence (Chi-Square $\chi^2$ test):} Fundamentally verifies that sequential spontaneous counting events appropriately and exclusively shadow probabilistic Poisson decay profiles unmolested by systematic electronic stutter, electrical spikes, or multi-pulsation PMT breakdown. Evaluating empirical dispersion against expected standard-error variations is formulated algebraically as:
    \begin{equation}
        \chi^2 = \sum_{i=1}^{n} \frac{(X_i - X_{av})^2}{X_{av}}
    \end{equation}
    where $n$ is sequential measurement repetition (e.g., 10 iterations), $X_i$ is net counts strictly per periodic repetition timeframe, and $X_{av}$ denotes the consolidated arithmetic mean established across all temporal iterations. Routine diagnostic confidence mandates the resultant $\chi^2$ aggregate score to fall securely within bounded upper-lower threshold tail distributions (0.05 \& 0.95 probability vectors) approximately $\geq 90\%$ of trials.

    \item \textbf{Minimum Detectable Activity (MDA):} Essentially calculates limits of detection. Taking a long temporal background measurement wholly devoid of extraneous local sources, the module assesses background noise standard deviation. It establishes absolute detectable quantity levels distinguishing statistical zero vs a fundamental positive count rate, mathematically defined closely by 3 standard deviations encompassing local room baseline averages.
\end{itemize}

\subsection{Radioactive Iodine Uptake (RAIU) Calculation Execution}
Obtaining the fractional retention quantitatively requires careful background variance differentiation across distinct counts representing identically matched specific doses. Measurements uniformly correspond mathematically to temporally defined spans (usually 2-hour early baseline versus 24-hour plateau assessments) utilizing standard calibration geometric metrics.
\begin{equation}
    \text{RAIU (\%)} = \left( \frac{P - T}{C - B} \right) \times 100
\end{equation}
where variables mathematically encapsulate structural conditions \cite{captus4000e}:
\begin{itemize}
    \item $P$ designates the gross integrated \textbf{Patient Neck count rate}, captured exclusively directly 25-30 cm anterior to the thyroid.
    \item $T$ denotes specific \textbf{Patient Background} (typically quantified by counts harvested from sweeping the thigh—structurally mirroring non-thyroidal somatic circulation scatter devoid of active trapping, ensuring vascular pooling counts are mathematically annulled).
    \item $C$ signifies \textbf{Standard Reference counts} initially derived via scanning the pre-administered identical matched tracer directly secured identically in the previously discussed specific lucite neck phantom geometrical housing.
    \item $B$ indicates pervasive intrinsic \textbf{Room Background counts}. 
\end{itemize}
These empirical corrections eliminate cosmic variations, vascular distributions, and scattered attenuation disparities. Extrapolating delayed fractions additionally relies functionally upon inherent decay-rate normalization adjustments depending heavily on software utilization or mathematical extrapolation algorithms matching initial ($T_0$) standard capacities equivalently mapped to observation delays (e.g., utilizing $A = A_0 e^{-\lambda t}$). Valid normal RAIU limits routinely exhibit characteristic 2-hour bounds between $5\% - 15\%$ tracking steadily up approaching $10\% - 35\%$ sequentially at the definitive 24-hour evaluation metric.
